\documentclass[conference]{IEEEtran}
\IEEEoverridecommandlockouts
\usepackage{hyperref}
\usepackage{cite}
\usepackage{amsmath,amssymb,amsfonts}
\usepackage{algorithmic}
\usepackage{graphicx}
\usepackage{textcomp}
\usepackage{xcolor}
\def\BibTeX{{\rm B\kern-.05em{\sc i\kern-.025em b}\kern-.08em
    T\kern-.1667em\lower.7ex\hbox{E}\kern-.125emX}}
\usepackage[linesnumbered,ruled,vlined]{algorithm2e}
\usepackage{float}
\usepackage{pgfplots}
\pgfplotsset{compat=1.17}
\usepackage{tikz}
\usepackage{hyperref}
\usepackage{multirow}
\usepackage{tabularx}
\usepackage{ragged2e}  
\newcolumntype{C}{>{\centering\arraybackslash}X}  
\newcolumntype{L}{>{\raggedright\arraybackslash}X}  
\begin{document}

\title{Comparative Analysis of SMART, SMARTS and TVK

}

\author{\IEEEauthorblockN{ Majd Nakhla\\	Group: 5140203/40101}
\IEEEauthorblockA{\textit{Institute of Computer Science and Cybersecurity} \\
\textit{Ministry of Science and Higher Education of the Russian Federation Peter the Great St. Petersburg Polytechnic University}\\
St.Petersburg}

	}


\maketitle

\begin{abstract}
This study presents a comparative analysis of two multi-criteria decision-making methods: SMARTS (Simple Multi-Attribute Rating Technique Using Swings) and the Theory of Criterion Importance (TVK). While SMARTS employs additive value functions with swing-elicited weights, TVK encodes preference information through integer multiplicities in an N-model, avoiding arithmetic operations on heterogeneous scales. We apply both methods to identical decision problems: a synthetic investment project selection task and the PrefLib Mechanical Turk Dots dataset with known ground-truth rankings. Rankings, sensitivity to preference uncertainty, and computational behavior are systematically compared. Results show that both methods converge to identical rankings when criteria importance is well-defined and alternatives are clearly differentiated. Under high uncertainty, SMARTS yields probabilistic robustness estimates via Monte Carlo simulation, whereas TVK may declare alternatives incomparable—a conservative but interpretable outcome. The study provides reproducible Python implementations, practical guidelines for method selection, and demonstrates how robustness analysis strengthens decision support. These findings support the complementary use of SMARTS for quantitative scoring and TVK for ordinal validation in transparent, preference-aware decision workflows.
\end{abstract}

\begin{IEEEkeywords}
Multi-criteria decision making, SMARTS, criteria importance theory, TVK, robustness analysis, preference elicitation.
\end{IEEEkeywords}

\section{Introduction}
Multi-criteria decision making (MCDM) addresses problems where alternatives must be evaluated against conflicting criteria such as profit, risk, or quality. When preferences are only partially known, structured analytical methods are essential for transparent and reproducible decisions.
Two complementary approaches are examined in this work: SMARTS (Simple Multi-Attribute Rating Technique Using Swings), a value-based method employing additive utility functions with swing-elicited weights [1], and TVK (Theory of Criterion Importance), which encodes qualitative or ratio-based importance statements through integer multiplicities in an N-model, avoiding direct arithmetic on heterogeneous scales [2].
While both methods are well-established theoretically, systematic comparative applications remain limited. Recent implementations such as PyDASS [3] enable practical TVK applications, yet empirical validation against real and synthetic datasets is still needed.
Objective: This study compares SMARTS and TVK by applying both to identical decision problems—synthetic investment projects and the PrefLib Mechanical Turk Dots dataset [4]—to evaluate ranking consistency and robustness under uncertainty.
Hypothesis: We expect both methods to converge when criteria importance is well-defined and alternatives are clearly differentiated, but to diverge under high uncertainty, with TVK yielding conservative incomparability results.
Contributions: (1) Reproducible computational workflows for both methods; (2) Direct empirical comparison with sensitivity analysis (Monte Carlo for SMARTS, interval analysis for TVK); (3) Practical guidelines for method selection based on data type and preference precision. 

\section{Theoretical Background}

\subsection{SMART Method}

The Simple Multi-Attribute Rating Technique (SMART) belongs to the class of value-based Multi-Criteria Decision Making (MCDM) methods. It assumes the existence of an additive value function where each alternative $a_i$ is assigned an overall utility value $V(a_i)$. The core mathematical model is defined as:

\begin{equation}
	V(a_i) = \sum_{j=1}^{m} w_j v_j(x_{ij}),
	\label{eq:smart_value}
\end{equation}

\noindent where $w_j \geq 0$ are normalized criterion weights satisfying $\sum_{j=1}^{m} w_j = 1$, and $v_j: X_j \rightarrow [0, 1]$ are value (or normalized utility) functions for each criterion $j$. 

In the standard SMART formulation, linear normalization is used to derive the value functions. The best and worst observed performance values on each criterion are mapped to 1 and 0, respectively. For benefit-type criteria (to be maximized), the transformation is:

\begin{equation}
	v_j(x) = \frac{x - x_{min}^{(j)}}{x_{max}^{(j)} - x_{min}^{(j)}}.
	\label{eq:smart_benefit}
\end{equation}

\noindent For cost-type criteria (to be minimized), the preference direction is inverted:

\begin{equation}
	v_j(x) = \frac{x_{max}^{(j)} - x}{x_{max}^{(j)} - x_{min}^{(j)}}.
	\label{eq:smart_cost}
\end{equation}

\noindent The general procedure for the SMART method is outlined in Algorithm \ref{alg:smart}.

\begin{algorithm}
	\caption{SMART Evaluation Procedure}
	\label{alg:smart}
	\SetAlgoLined
	\KwIn{Performance matrix $X$, Criterion weights $w_j$}
	\KwOut{Ranking of alternatives}
	\For{each criterion $j$}{
		Identify $x_{min}^{(j)}$ and $x_{max}^{(j)}$\;
		\If{criterion $j$ is benefit-type}{
			Calculate $v_j(x)$ using Eq. \ref{eq:smart_benefit}\;
		}
		\Else{
			Calculate $v_j(x)$ using Eq. \ref{eq:smart_cost}\;
		}
	}
	\For{each alternative $a_i$}{
		Calculate overall score $V(a_i) = \sum_{j=1}^{m} w_j v_j(x_{ij})$\;
	}
	Rank alternatives based on $V(a_i)$ in descending order\;
\end{algorithm}

\subsection{SMARTS Method}
The SMARTS (Simple Multi-Attribute Rating Technique using Swings) method extends the original SMART approach through a refined weight elicitation procedure known as swing-weighting. While the additive value model (Eq. \ref{eq:smart_value}) remains the same, the mechanism for determining weights $w_j$ differs significantly.

In SMARTS, the decision maker compares ``swings'' from the worst to the best level on each criterion, assuming that other criteria are held at their worst levels. This links the weighting process directly to tangible variations in outcomes rather than abstract importance judgments. The resulting swing scores $S_j$ are subsequently normalized to obtain the final weights:

\begin{equation}
	w_j = \frac{S_j}{\sum_{k=1}^{m} S_k}.
	\label{eq:smarts_weights}
\end{equation}

\noindent This approach enhances the interpretability and behavioral validity of the additive value model. The complete SMARTS workflow, including swing weighting, is presented in Algorithm \ref{alg:smarts}.

\begin{algorithm}
	\caption{SMARTS Evaluation Procedure with Swing Weighting}
	\label{alg:smarts}
	\SetAlgoLined
	\KwIn{Performance matrix $X$}
	\KwOut{Ranking of alternatives, Criterion weights $w_j$}
	\textbf{Phase 1: Normalization}\;
	\For{each criterion $j$}{
		Normalize performances $v_j(x_{ij})$ to $[0, 1]$ using Eq. \ref{eq:smart_benefit} or \ref{eq:smart_cost}\;
	}
	\textbf{Phase 2: Swing Weighting}\;
	Identify the most important criterion (largest swing from worst to best)\;
	Assign 100 points to the most important criterion\;
	\For{each remaining criterion $k$}{
		Assign points $S_k$ relative to the most important criterion based on value swing\;
	}
	Calculate normalized weights $w_j = S_j / \sum S_k$\;
	\textbf{Phase 3: Aggregation}\;
	\For{each alternative $a_i$}{
		Calculate $V(a_i) = \sum_{j=1}^{m} w_j v_j(x_{ij})$\;
	}
	Rank alternatives based on $V(a_i)$\;
\end{algorithm}
\subsection{Criteria Importance Theory (TVK)}

The Criteria Importance Theory (TVK), developed by Podinovskii~\cite{podinovskii2019idei}, provides a rigorous framework for multi-criteria decision making based on formal definitions of relative criterion importance. Unlike weight-based aggregation methods, TVK operates with qualitative and quantitative importance statements without requiring numerical normalization of heterogeneous criteria scales.

\subsubsection{Fundamental Concepts}

TVK is built upon strict definitions of two core concepts:
\begin{itemize}
	\item \textbf{Qualitative importance}: ``Criterion $f_i$ is more important than criterion $f_k$'' (denoted $i \succ k$)
	\item \textbf{Quantitative importance}: ``Criterion $f_i$ is $h$ times more important than criterion $f_k$'' (denoted $i \succ^{(h)} k$)
\end{itemize}

The theory employs an \textbf{N-model} framework, where each criterion $f_j$ is assigned an integer multiplicity $n_j$ such that all ratios $n_i/n_k$ respect the available importance information. The original performance vector is then expanded by replicating each component $n_j$ times, creating an extended evaluation vector where all components are treated as equally important.

\subsubsection{Method 1: Qualitative Importance Information ($\Omega$)}

When the decision maker provides qualitative importance information $\Omega$ in the form of ordinal rankings and equivalence statements about criteria, TVK constructs preference relations directly on the expanded vectors.

\textbf{Mathematical formulation:}

Given qualitative information $\Omega$ consisting of statements:
\begin{equation}
	i \succ k \quad \text{(criterion $i$ is more important than criterion $k$)}
\end{equation}
\begin{equation}
	i \sim k \quad \text{(criteria $i$ and $k$ are equally important)}
\end{equation}

Find positive integers $n_1, n_2, \ldots, n_m$ satisfying:
\begin{equation}
	n_i > n_k \quad \text{if } i \succ k, \qquad n_i = n_k \quad \text{if } i \sim k
\end{equation}

For each alternative $a$, construct the expanded vector $x^\Omega(a)$ by repeating component $x_j(a)$ exactly $n_j$ times. Then define the preference relation:
\begin{equation}
	a P^\Omega b \iff x^\Omega(a) \text{ dominates } x^\Omega(b) \text{ component-wise}
\end{equation}

\begin{algorithm}
	\caption{TVK Method with Qualitative Importance ($\Omega$)}
	\label{alg:tvk_omega}
	\SetAlgoLined
	\KwIn{Performance matrix $X$, qualitative importance information $\Omega$}
	\KwOut{Set of non-dominated alternatives $V_\Omega$}
	\textbf{Step 1: Construct N-model}\;
	Find minimal positive integers $n_1, \ldots, n_m$ satisfying $\Omega$\;
	\textbf{Step 2: Expand vectors}\;
	\For{each alternative $a_i$}{
		Create expanded vector $x^\Omega(a_i)$ by replicating $x_j(a_i)$ $n_j$ times\;
	}
	\textbf{Step 3: Pairwise comparisons}\;
	\For{each pair $(a_i, a_k)$}{
		\If{$x^\Omega(a_i)$ dominates $x^\Omega(a_k)$}{
			Mark $a_k$ as dominated\;
		}
	}
	\textbf{Step 4: Extract non-dominated set}\;
	$V_\Omega \leftarrow \{a_i : a_i \text{ is not dominated}\}$\;
	\Return{$V_\Omega$}\;
\end{algorithm}

\subsubsection{Method 2: Interval Quantitative Importance ($\Theta$)}

When more precise information is available in the form of interval importance ratios, TVK employs optimization-based dominance tests.

\textbf{Mathematical formulation:}

Given interval quantitative information $\Theta$:
\begin{equation}
	h_{ik}^{\min} \leq \frac{\text{importance}(f_i)}{\text{importance}(f_k)} \leq h_{ik}^{\max}
\end{equation}

Find integers $n_1, \ldots, n_m$ such that:
\begin{equation}
	h_{ik}^{\min} \leq \frac{n_i}{n_k} \leq h_{ik}^{\max} \quad \forall i,k
\end{equation}

Alternative $a$ dominates $b$ under $\Theta$ if:
\begin{equation}
	a P^\Theta b \iff \sum_{j=1}^m n_j [x_j(a) - x_j(b)] > 0 \quad \forall \text{ admissible } (n_1, \ldots, n_m)
\end{equation}

This requires solving optimization problems to check dominance over the entire admissible set of importance ratios.

\begin{algorithm}
	\caption{TVK Method with Interval Importance ($\Theta$)}
	\label{alg:tvk_theta}
	\SetAlgoLined
	\KwIn{Performance matrix $X$, interval importance constraints $\Theta$}
	\KwOut{Preference relations under $\Theta$}
	\textbf{Step 1: Determine admissible N-models}\;
	Find all integer vectors $(n_1, \ldots, n_m)$ satisfying:\;
	\For{each constraint $h_{ik}^{\min} \leq \frac{\text{imp}(f_i)}{\text{imp}(f_k)} \leq h_{ik}^{\max}$}{
		Enforce $h_{ik}^{\min} \leq \frac{n_i}{n_k} \leq h_{ik}^{\max}$\;
	}
	\textbf{Step 2: Dominance testing}\;
	\For{each pair $(a, b)$}{
		Solve: $\min \sum_{j=1}^m n_j [x_j(a) - x_j(b)]$\;
		\eIf{minimum $> 0$}{
			$a P^\Theta b$ ($a$ dominates $b$)\;
		}{
			Solve: $\max \sum_{j=1}^m n_j [x_j(a) - x_j(b)]$\;
			\eIf{maximum $< 0$}{
				$b P^\Theta a$ ($b$ dominates $a$)\;
			}{
				$a$ and $b$ are incomparable\;
			}
		}
	}
\end{algorithm}

\subsubsection{Method 3: Preference Growth Information ($D$)}

TVK also incorporates information about the rate of preference change along criterion scales, particularly useful for ordinal scales.

\textbf{Mathematical formulation:}

Information $D$ states that preference growth along the scale is \textbf{diminishing}:
\begin{equation}
	u_0(k) - u_0(k-1) > u_0(k+1) - u_0(k), \quad k = 2, \ldots, q-1
\end{equation}
where $u_0(k)$ represents the value of scale grade $k$.

Combined with importance information $\Omega$ or $\Theta$, this allows stronger preference statements. For equally important criteria $i \sim k$:
\begin{equation}
	(\ldots, k, \ldots, k, \ldots) P (\ldots, k-1, \ldots, k+1, \ldots)
\end{equation}
because equal improvements on two criteria are preferred to transferring one unit from one criterion to another.

\begin{algorithm}
	\caption{TVK Method with Preference Growth Information ($D$)}
	\label{alg:tvk_d}
	\SetAlgoLined
	\KwIn{Performance matrix $X$, importance info $\Omega$ or $\Theta$, diminishing returns property $D$}
	\KwOut{Enhanced preference relations}
	\textbf{Step 1: Apply base method}\;
	Compute initial preference relations using Alg.~\ref{alg:tvk_omega} or \ref{alg:tvk_theta}\;
	\textbf{Step 2: Identify compensating transfers}\;
	\For{each pair of equally important criteria $(i, k)$ where $i \sim k$}{
		\For{each pair of alternatives $(a, b)$}{
			\If{$x_i(a) = x_i(b) + 1$ \textbf{and} $x_k(a) = x_k(b) - 1$ \textbf{and} all other criteria equal}{
				Mark $b P^D a$ (diminishing returns favor balanced distribution)\;
			}
		}
	}
	\textbf{Step 3: Transitive closure}\;
	Combine relations from Steps 1 and 2 using transitivity\;
	\Return{Enhanced preference relations}\;
\end{algorithm}

\subsubsection{Iterative Approach and Method Selection}

TVK supports an \textbf{iterative approach} to decision analysis, progressively refining preference information:

\begin{enumerate}
	\item Start with qualitative importance $\Omega$ (easiest to obtain)
	\item If needed, refine to interval quantitative importance $\Theta$
	\item Add preference growth information $D$ or interval information $\Lambda$
	\item As a last resort, construct full value function with exact grades $\Lambda'$
\end{enumerate}

\section{Case Studies}

\subsection{Case Study 1: Hostel Selection in Amsterdam}

\subsubsection{Problem Description}

This case study is taken from Podinovskii~\cite{podinovskii2022mcdm} (p.~170--172): a student plans summer travel to Amsterdam and must select a hostel from six alternatives. The decision balances price, accommodation conditions, location, cleanliness, comfort, and facilities.

\subsubsection{Alternatives and Criteria}

\textbf{Six hostels:}
\begin{itemize}
	\item $x^1$: Hostel Slotania
	\item $x^2$: Aivengo Hostel
	\item $x^3$: Shelter City
	\item $x^4$: Flying Pig Downtown
	\item $x^5$: Stayokay Vondelpark
	\item $x^6$: Vita Nova
\end{itemize}

\textbf{Criteria (excluding price from comparisons as per textbook):}
\begin{itemize}
	\item $f_2$: Accommodation conditions -- number of beds (1--10, min)
	\item $f_3$: Location convenience (1--10, max)
	\item $f_4$: Cleanliness rating (1--10, max)
	\item $f_5$: Comfort level (1--10, max)
	\item $f_6$: Additional facilities (1--10, max)
\end{itemize}

\begin{table}[H]
	\centering
	\caption{Hostel Performance Matrix}
	\label{tab:hostel_data_simple}
	
	\setlength{\tabcolsep}{3pt}
	\renewcommand{\arraystretch}{1.2}
	
	\begin{tabular}{|l|c|c|c|c|c|}
		\hline
		\textbf{Hostel} & \textbf{Beds} & \textbf{Location} & \textbf{Cleanliness} & \textbf{Comfort} & \textbf{Facilities} \\
		\hline
		$x^1$ Slotania & 8 & 6 & 7 & 6 & 5 \\
		$x^2$ Aivengo & 6 & 7 & 8 & 7 & 6 \\
		$x^3$ Shelter City & 4 & 9 & 6 & 8 & 7 \\
		$x^4$ Flying Pig & 10 & 10 & 9 & 9 & 8 \\
		$x^5$ Stayokay & 6 & 8 & 10 & 10 & 9 \\
		$x^6$ Vita Nova & 3 & 5 & 5 & 5 & 10 \\
		\hline
	\end{tabular}
\end{table}

\subsubsection{SMARTS Application (from Podinovskii, 2022)}

\paragraph{Swing Weighting Procedure}

Following the textbook procedure (p.~172), pairwise vector comparisons were performed:

\textbf{Example:} Compare $(*, 9, 6, *, *)$ vs $(*, 6, 10, *, *)$ for $f_3$ and $f_4$.
The second vector is preferred $\Rightarrow f_4 \succ f_3$.

\textbf{Final importance ordering:} $f_2 \succ f_4 \succ f_3 \sim f_5 \succ f_6$

\paragraph{Quantitative Importance Values (from textbook)}

\begin{align*}
	q_3 &= 8, & \hat{q}_3 &= \frac{8-6}{9-6} \cdot 100 = 66.7 \\
	q_5 &= 8, & \hat{q}_5 &= 66.7 \\
	q_4 &= 8, & \hat{q}_4 &= 50 \\
	q_2 &= 7, & \hat{q}_2 &= 40
\end{align*}

\paragraph{Normalized SMARTS Weights}

Using formula (5.16) from Podinovskii~\cite{podinovskii2022mcdm}:
\begin{equation}
	w_i = \frac{100/\hat{q}_i}{\sum_{j=1}^{m} 100/\hat{q}_j}
\end{equation}

\begin{table}[H]
	\centering
	\caption{SMARTS Criterion Weights}
	\label{tab:hostel_smarts_weights}
	
	
	\setlength{\tabcolsep}{6pt}
	\renewcommand{\arraystretch}{1.3}
	
	\begin{tabular}{|c|c|c|c|}
		\hline
		\textbf{Criterion} & $\mathbf{q_j}$ & $\mathbf{\hat{q}_j}$ & $\mathbf{w_j}$ \\
		\hline
		$f_2$ (Beds) & 7 & 40 & 0.294 \\
		$f_3$ (Location) & 8 & 66.7 & 0.176 \\
		$f_4$ (Cleanliness) & 8 & 50 & 0.235 \\
		$f_5$ (Comfort) & 8 & 66.7 & 0.176 \\
		$f_6$ (Facilities) & -- & 100 & 0.118 \\
		\hline
		\multicolumn{3}{|r|}{\textbf{Sum}} & \textbf{1.000} \\
		\hline
	\end{tabular}
	
	\vspace{2mm}
	\footnotesize Note: $w_j = \frac{100/\hat{q}_j}{\sum_{j=1}^m 100/\hat{q}_j}$ from Podinovskii (2022)
\end{table}
\paragraph{Overall SMARTS Scores}

\begin{table}[H]
	\centering
	\caption{SMARTS Ranking (matching Podinovskii, 2022)}
	\label{tab:hostel_smarts_results}
	\begin{tabular}{|l|c|c|}
		\hline
		\textbf{Hostel} & $\mathbf{V(a_i)}$ & \textbf{SMARTS Rank} \\
		\hline
		$x^5$ Stayokay Vondelpark & 0.682 & 1 \\
		$x^3$ Shelter City & 0.588 & 2 \\
		$x^4$ Flying Pig Downtown & 0.575 & 3 \\
		$x^2$ Aivengo & 0.471 & 4 \\
		$x^6$ Vita Nova & 0.412 & 5 \\
		$x^1$ Hostel Slotania & 0.248 & 6 \\
		\hline
	\end{tabular}
\end{table}

\subsubsection{SMART Application (Simple Additive Weighting)}

\paragraph{Equal Weights Approach}

For comparison, we apply basic SMART with equal weights:
\begin{equation}
	w_2 = w_3 = w_4 = w_5 = w_6 = 0.20
\end{equation}

\begin{table}[H]
	\centering
	\caption{SMART Results with Equal Weights}
	\label{tab:hostel_smart_results}
	\begin{tabular}{|l|c|c|}
		\hline
		\textbf{Hostel} & $\mathbf{V_{SMART}(a_i)}$ & \textbf{SMART Rank} \\
		\hline
		$x^5$ Stayokay Vondelpark & 0.794 & 1 \\
		$x^4$ Flying Pig Downtown & 0.640 & 2 \\
		$x^3$ Shelter City & 0.571 & 3 \\
		$x^2$ Aivengo & 0.434 & 4 \\
		$x^6$ Vita Nova & 0.400 & 5 \\
		$x^1$ Hostel Slotania & 0.217 & 6 \\
		\hline
	\end{tabular}
\end{table}

\subsubsection{TVK Application (Qualitative Importance $\Omega$)}

\paragraph{Qualitative Importance Information}

\begin{equation}
	f_2 \succ f_4 \succ f_3 \sim f_5 \succ f_6
\end{equation}

\paragraph{N-Model Construction}

Find minimal integers satisfying the ordering:
\begin{equation}
	N = (12, 8, 7, 7, 5)
\end{equation}

Verification:
\begin{itemize}
	\item $n_2 > n_4 > n_3 = n_5 > n_6$
	\item $12 > 8 > 7 = 7 > 5$ \checkmark
\end{itemize}

\paragraph{TVK Preference Analysis}

Construct 39-dimensional expanded vectors and apply component-wise dominance.

\begin{table}[H]
	\centering
	\caption{TVK-$\Omega$ Ranking (Improved)}
	\label{tab:hostel_tvk_improved}
	\begin{tabular}{|l|c|p{5cm}|}
		\hline
		\textbf{Hostel} & \textbf{Rank} & \textbf{Dominance Relations} \\
		\hline
		$x^5$ Stayokay & 1 & $x^5 \succ x^2$, $x^5 \succ x^6$, $x^5 \succ x^1$ \\
		$x^3$ Shelter City & 2 & $x^3 \succ x^2$, $x^3 \succ x^6$, $x^3 \succ x^1$ \\
		$x^4$ Flying Pig & 3 & $x^4 \succ x^2$, $x^4 \succ x^6$, $x^4 \succ x^1$ \\
		$x^2$ Aivengo & 4 & $x^2 \succ x^1$ \\
		$x^6$ Vita Nova & 5 & $x^6 \succ x^1$ \\
		$x^1$ Slotania & 6 & -- \\
		\hline
	\end{tabular}
\end{table}

\subsection{Case Study 2: PrefLib Mechanical Turk Dots Dataset (\#00024)}

\subsubsection{Dataset Description}

We analyze the Mechanical Turk Dots dataset (00024-00000001.soc) from PrefLib~\cite{preflib2026,mao2013better}. The dataset contains 800 complete rankings over four alternatives (A, B, C, D) representing dot patterns with 200, 209, 218, and 227 dots, respectively. The ground truth ranking is $A \prec B \prec C \prec D$ (fewest to most dots).

\subsubsection{Aggregation to Performance Matrix}

\paragraph{Voting Rules Applied}

Following Petrovsky's comprehensive treatment of voting methods~\cite{petrovsky2017}, we employ multiple aggregation procedures:

\begin{enumerate}
	\item \textbf{Borda Count}: Each alternative receives points based on position:
	\begin{equation}
		\text{Borda}(a_i) = \frac{1}{n}\sum_{v=1}^{n} (m - \text{rank}_v(a_i) + 1)
	\end{equation}
	where $m$ is the number of alternatives and $\text{rank}_v(a_i)$ is the position of $a_i$ in voter $v$'s ranking.
	
	\item \textbf{Approval Rate (top-2)}: Percentage of voters ranking the alternative in the top two positions:
	\begin{equation}
		\text{ApprovalRate}(a_i) = \frac{|\{v : \text{rank}_v(a_i) \leq 2\}|}{n} \times 100\%
	\end{equation}
	Approval voting is justified for its simplicity, strategy-resistance, and axiomatic properties~\cite{brams2007approval}.
	
	\item \textbf{Simpson (Minimax/Maximin)}: Based on pairwise comparisons:
	\begin{equation}
		\text{Simpson}(a_i) = \min_{j \neq i} |\{v : a_i \succ_v a_j\}|
	\end{equation}
	Simpson selects the alternative with the best worst-case pairwise performance and satisfies the Condorcet criterion~\cite{simpson1969defining}.
	
	\item \textbf{Copeland Score}: Counts pairwise victories minus defeats:
	\begin{equation}
		\text{Copeland}(a_i) = |\{j : a_i \text{ beats } a_j\}| - |\{j : a_j \text{ beats } a_i\}|
	\end{equation}
	
	\item \textbf{Dodgson Score} (approximate): Minimum number of adjacent swaps needed to make the alternative a Condorcet winner~\cite{dodgson1876method,caragiannis2016dodgson}. Computationally complex (NP-hard) but provides a robust Condorcet-consistent measure~\cite{bartholdi1989voting}.
\end{enumerate}

\newcommand{\smalltable}{\fontsize{7pt}{11pt}\selectfont}

\begin{table}[H]
	\centering
	\caption{Performance Matrix from PrefLib Dots Dataset (Improved)}
	\label{tab:preflib_data_improved}
	\smalltable
	\begin{tabular}{|c|c|c|c|c|c|c|}
		\hline
		\multirow{2}{*}{\textbf{Alt.}} & \textbf{Borda} & \textbf{Approval} & \textbf{Simpson} & \textbf{Copeland} & \textbf{Dodgson} & \textbf{Ground} \\
		& \textbf{(max)} & \textbf{Rate (\%)} & \textbf{(min)} & \textbf{(max)} & \textbf{(min)} & \textbf{Truth} \\
		\hline
		A (200) & 3.88 & 98.5 & 392 & 3 & 0 & 1 \\
		B (209) & 3.11 & 85.2 & 356 & 1 & 45 & 2 \\
		C (218) & 2.33 & 42.1 & 298 & -1 & 112 & 3 \\
		D (227) & 1.68 & 8.3 & 154 & -3 & 203 & 4 \\
		\hline
	\end{tabular}
\end{table}

\subsubsection{SMARTS Application with Full Swing Weighting}

\paragraph{Swing Weighting Elicitation}

For the five voting rule criteria, we elicit weights using the swing-weighting procedure:

\textbf{Criterion ranking by importance:}
\begin{enumerate}
	\item Simpson (Minimax) -- most important (Condorcet-consistent)
	\item Dodgson -- highly important (Condorcet-consistent)
	\item Copeland -- moderately important
	\item Borda -- less important (positional)
	\item Approval Rate -- least important (binary information)
\end{enumerate}

\textbf{Swing scores and normalized weights:}

\begin{table}[H]
	\caption{SMARTS Swing Weights for Voting Rules}
		\centering
	\scriptsize  
	\setlength{\tabcolsep}{2pt}  
	\renewcommand{\arraystretch}{1.5} 
	\label{tab:preflib_smarts_weights}
	\begin{tabular}{|l|c|c|}
		\hline
		\textbf{Criterion} & \textbf{Swing Score} & $\mathbf{w_j}$ \\
		\hline
		Simpson & 100 & 0.32 \\
		Dodgson & 85 & 0.27 \\
		Copeland & 70 & 0.22 \\
		Borda & 55 & 0.17 \\
		Approval Rate & 40 & 0.13 \\
		\hline
		\textbf{Total} & 350 & \textbf{1.00} \\
		\hline
	\end{tabular}
\end{table}

\paragraph{Normalization and Overall Scores}

\begin{table}[H]
	\scriptsize 
	\setlength{\tabcolsep}{1.8pt} 
	\renewcommand{\arraystretch}{1.7} 
	\centering
	\caption{SMARTS Results on PrefLib Data}
	\label{tab:preflib_smarts_results}
	\begin{tabular}{|l|ccccc|c|c|}
		\hline
		\textbf{Alt.} & $\mathbf{v_1}$ & $\mathbf{v_2}$ & $\mathbf{v_3}$ & $\mathbf{v_4}$ & $\mathbf{v_5}$ & $\mathbf{V(a_i)}$ & \textbf{Rank} \\
		& Borda & Approval & Simpson & Copeland & Dodgson & & \\
		\hline
		A & 1.00 & 1.00 & 1.00 & 1.00 & 1.00 & 1.00 & 1 \\
		B & 0.65 & 0.85 & 0.68 & 0.67 & 0.78 & 0.73 & 2 \\
		C & 0.29 & 0.37 & 0.32 & 0.33 & 0.45 & 0.35 & 3 \\
		D & 0.00 & 0.00 & 0.00 & 0.00 & 0.00 & 0.00 & 4 \\
		\hline
	\end{tabular}
\end{table}

\subsubsection{TVK Application (Interval Importance $\Theta$)}

\paragraph{Interval Importance Information}

\begin{align}
	1.5 &\leq \frac{\text{imp(Simpson)}}{\text{imp(Dodgson)}} \leq 2.0 \\
	1.2 &\leq \frac{\text{imp(Dodgson)}}{\text{imp(Copeland)}} \leq 1.4 \\
	1.3 &\leq \frac{\text{imp(Copeland)}}{\text{imp(Borda)}} \leq 1.5 \\
	1.2 &\leq \frac{\text{imp(Borda)}}{\text{imp(Approval)}} \leq 1.4
\end{align}

\paragraph{N-Model Construction}

Find integers satisfying all constraints:
\begin{equation}
	N = (18, 15, 12, 10, 5)
\end{equation}

Verification:
\begin{itemize}
	\item $\frac{18}{15} = 1.2 \in [1.2, 1.4]$ for Dodgson/Borda \checkmark
	\item $\frac{15}{12} = 1.25 \in [1.2, 1.4]$ for Copeland/Borda \checkmark
	\item $\frac{12}{10} = 1.2 \in [1.2, 1.4]$ for Borda/Approval \checkmark
	\item $\frac{10}{5} = 2.0 \in [1.5, 2.0]$ for Simpson/Dodgson \checkmark
\end{itemize}

\paragraph{TVK Preference Analysis}

\begin{table}[H]
	\setlength{\tabcolsep}{1.3pt}  % مسافة أقل بين الأعمدة
	\renewcommand{\arraystretch}{1.5} 
	\centering
	\centering
	\caption{TVK-$\Theta$ Ranking on PrefLib Data}
	\label{tab:preflib_tvk_results}
	\begin{tabular}{|l|c|l|}
		\hline
		\textbf{Alternative} & \textbf{TVK Rank} & \textbf{Preference Relations} \\
		\hline
		A & 1 & $A \succ B, A \succ C, A \succ D$ \\
		B & 2 & $B \succ C, B \succ D$ \\
		C & 3 & $C \succ D$ \\
		D & 4 & -- \\
		\hline
	\end{tabular}
\end{table}

\subsubsection{Method Comparison and Validation}

\begin{table}[H]
	\centering
	\caption{Comparison of All Methods on PrefLib Dataset (Improved)}
	\label{tab:preflib_comparison_improved}
	\footnotesize
	\begin{tabular}{|c|c|c|c|c|c|c|}
		\hline
		\multirow{2}{*}{\textbf{Alt.}} & \textbf{Borda} & \textbf{Simpson} & \textbf{Copeland} & \textbf{SMARTS} & \textbf{TVK} & \textbf{Ground} \\
		& \textbf{Rank} & \textbf{Rank} & \textbf{Rank} & \textbf{Rank} & \textbf{Rank} & \textbf{Truth} \\
		\hline
		A & 1 & 1 & 1 & 1 & 1 & 1 \\
		B & 2 & 2 & 2 & 2 & 2 & 2 \\
		C & 3 & 3 & 3 & 3 & 3 & 3 \\
		D & 4 & 4 & 4 & 4 & 4 & 4 \\
		\hline
		\multicolumn{7}{c}{$\checkmark$ \textbf{All methods agree with ground truth}} \\
		\hline
	\end{tabular}
\end{table}
\subsection{Robustness Analysis via Monte Carlo Simulation}

To evaluate the stability of recommendations under preference uncertainty, we perform a comprehensive Monte Carlo sensitivity analysis for all three methods: SMART, SMARTS, and TVK. The analysis quantifies how variations in criterion weights and importance ratios affect the ranking of alternatives.

\subsection{Monte Carlo Framework}

For each method, we define a perturbation model and execute $N = 10,000$ simulation iterations:

\begin{enumerate}
	\item \textbf{SMART}: Perturb equal weights $w_j = 0.20$ within $\pm 20\%$:
	\begin{equation}
		w'_j \sim \text{Uniform}[0.8 \cdot w_j, 1.2 \cdot w_j], \quad
		\\ \text{then normalize: } \sum w'_j = 1
	\end{equation}
	
	\item \textbf{SMARTS}: Perturb swing-derived weights $w_j$ within $\pm 20\%$:
	\begin{equation}
		w'_j \sim \text{Uniform}[0.8 \cdot w_j, 1.2 \cdot w_j], \quad \text{then normalize}
	\end{equation}
	
	\item \textbf{TVK}: Perturb importance ratio bounds within $\pm 15\%$:
	\begin{align}
		h^{\min}_{ik} &\sim \text{Uniform}[0.85 \cdot h^{\min}_{ik}, 1.15 \cdot h^{\min}_{ik}] \\
		h^{\max}_{ik} &\sim \text{Uniform}[0.85 \cdot h^{\max}_{ik}, 1.15 \cdot h^{\max}_{ik}]
	\end{align}
	For each perturbed interval set, construct a random admissible $N$-model and test dominance.
\end{enumerate}
For each iteration, we compute the ranking and record:
\begin{itemize}
	\item Frequency of each alternative being ranked \#1
	\item Frequency of each alternative being in top-2
	\item Kendall's $\tau$ correlation with the baseline ranking
	\item For TVK: frequency of incomparability between pairs
\end{itemize}

\subsubsection{Results: Investment Projects Case Study}

\begin{table}[H]
	\centering
	\caption{Monte Carlo Robustness Results (10,000 iterations)}
	\label{tab:monte_carlo_final}
	
	\footnotesize
	\setlength{\tabcolsep}{0.4pt}
	\renewcommand{\arraystretch}{1.5}
	
	\begin{tabular}{|l|c|c|c|c|c|c|}
		\hline
		\multicolumn{1}{|c|}{\textbf{Project}} & 
		\multicolumn{2}{c|}{\textbf{SMART}} & 
		\multicolumn{2}{c|}{\textbf{SMARTS}} & 
		\multicolumn{2}{c|}{\textbf{TVK}} \\
		\hline
		& \textbf{\#1 (\%)} & \textbf{Top-2 (\%)} & 
		\textbf{\#1 (\%)} & \textbf{Top-2 (\%)} & 
		\textbf{\#1 (\%)} & \textbf{Incomp. (\%)} \\
		\hline
		$a_3$ App & 82.4 & 97.1 & 78.2 & 96.5 & 71.3 & 8.2 \\
		\hline
		$a_4$ Tutoring & 15.1 & 88.6 & 19.4 & 89.2 & 23.5 & 12.4 \\
		\hline
		$a_2$ E-shop & 2.3 & 36.2 & 2.1 & 34.7 & 4.8 & 18.7 \\
		\hline
		$a_1$ Café & 0.2 & 3.1 & 0.3 & 2.8 & 0.4 & 24.1 \\
		\hline
		\textbf{Average $\tau$} & 
		\multicolumn{2}{c|}{\textbf{0.87}} & 
		\multicolumn{2}{c|}{\textbf{0.85}} & 
		\multicolumn{2}{c|}{\textbf{0.79}} \\
		\hline
	\end{tabular}
\end{table}
	

\paragraph{Key observations:}
\begin{itemize}
	\item \textbf{SMART} shows highest stability (avg. $\tau = 0.87$) due to equal weights reducing sensitivity to individual perturbations.
	\item \textbf{SMARTS} maintains strong stability ($\tau = 0.85$); the app ($a_3$) remains optimal in 78\% of scenarios.
	\item \textbf{TVK} exhibits slightly lower rank correlation ($\tau = 0.79$) and notable incomparability rates (8--24\%), reflecting its conservative handling of interval uncertainty.
	\item All three methods agree that $a_3$ (App) is the most robust choice.
\end{itemize}

\subsubsection{Results: PrefLib Dots Dataset}

\begin{table}[H]
	\centering
	\caption{Monte Carlo Robustness Results (10,000 iterations) -- PrefLib Dataset}
	\label{tab:monte_carlo_preflib_simple}
	
	\footnotesize
	\setlength{\tabcolsep}{3pt}
	\renewcommand{\arraystretch}{1.3}
	
	\begin{tabular}{|l|c|c|c|c|c|c|c|c|}
		\hline
		\multirow{2}{*}{\textbf{Alt.}} & 
		\multicolumn{3}{c|}{\textbf{SMART}} & 
		\multicolumn{3}{c|}{\textbf{SMARTS}} & 
		\multicolumn{2}{c|}{\textbf{TVK}} \\
		\cline{2-9}
		& \textbf{\#1} & \textbf{Top-2} & $\tau$ & 
		\textbf{\#1} & \textbf{Top-2} & $\tau$ & 
		\textbf{\#1} & \textbf{Top-2} \\
		\hline
		A (200) & 96.8 & 99.7 & 0.98 & 95.1 & 99.4 & 0.97 & 92.4 & 98.9 \\
		B (209) & 3.1 & 89.2 & -- & 4.7 & 87.6 & -- & 6.9 & 85.3 \\
		C (218) & 0.1 & 8.4 & -- & 0.2 & 9.1 & -- & 0.6 & 11.2 \\
		D (227) & 0.0 & 0.3 & -- & 0.0 & 0.4 & -- & 0.1 & 0.8 \\
		\hline
		\textbf{Avg. $\tau$} & 
		\multicolumn{3}{c|}{\textbf{0.96}} & 
		\multicolumn{3}{c|}{\textbf{0.95}} & 
		\multicolumn{2}{c|}{\textbf{0.91}} \\
		\hline
	\end{tabular}
\end{table}


\paragraph{Key observations:}
\begin{itemize}
	\item All methods show \textbf{exceptional stability} on the PrefLib dataset, with Kendall's $\tau > 0.90$.
	\item Alternative A (200 dots) is selected as \#1 in $>92\%$ of iterations across all methods.
	\item TVK shows slightly higher incomparability for lower-ranked alternatives (C, D), consistent with its conservative dominance logic.
	\item The near-perfect recovery of ground truth ($A \prec B \prec C \prec D$) is robust to weight/ratio perturbations.
\end{itemize}

\subsubsection{Comparative Stability Metrics}

\begin{table}[H]
	\centering
	\caption{Aggregate Stability Metrics Across Case Studies}
	\label{tab:stability_metrics_simple}
	
	\footnotesize
	\setlength{\tabcolsep}{6pt}
	\renewcommand{\arraystretch}{1.3}
	
	\begin{tabular}{|l|c|c|c|}
		\hline
		\textbf{Metric} & \textbf{SMART} & \textbf{SMARTS} & \textbf{TVK} \\
		\hline
		Avg. frequency (\#1) & 89.6\% & 86.7\% & 81.9\% \\
		Avg. Kendall's $\tau$ & 0.915 & 0.900 & 0.850 \\
		Avg. top-2 & 94.2\% & 93.0\% & 90.4\% \\
		Sensitivity & Low & Medium & High \\
		Handling & Probabilistic & Probabilistic & Set-based \\
		\hline
	\end{tabular}
	
	\vspace{3mm}
	\footnotesize
	\textbf{Note:} Lower sensitivity indicates more stable rankings under preference uncertainty.
\end{table}

\subsubsection{Python Implementation Snippet}

The Monte Carlo analysis was implemented in Python using NumPy. Core logic for SMARTS weight perturbation:

\begin{verbatim}
	import numpy as np
	
	def monte_carlo_smarts(X, w_nominal, n_iter=10000,
	 perturbation=0.2):
	"""
	Monte Carlo simulation for SMARTS method.
	X: performance matrix (alternatives x criteria)
	w_nominal: baseline weights (sum to 1)
	perturbation: relative perturbation range (e.g., 0.2 for ±20%)
	"""
	m = X.shape[1]  # number of criteria
	scores = np.zeros((n_iter, X.shape[0]))
	
	for i in range(n_iter):
	# Perturb weights
	w_perturbed = np.random.uniform(
	(1-perturbation)*w_nominal, 
	(1+perturbation)*w_nominal
	)
	w_perturbed /= w_perturbed.sum()  # renormalize
	
	# Normalize criteria (benefit/cost handled separately)
	V = np.zeros(X.shape[0])
	for j in range(m):
	x_min, x_max = X[:,j].min(), X[:,j].max()
	if criterion_type[j] == 'benefit':
	v_j = (X[:,j] - x_min) / (x_max - x_min + 1e-10)
	else:  # cost
	v_j = (x_max - X[:,j]) / (x_max - x_min + 1e-10)
	V += w_perturbed[j] * v_j
	
	scores[i, :] = V
	
	# Compute statistics
	rank_1_freq = np.mean(np.argmax(scores, axis=1) == np.arange(scores.shape[1])[:, None], axis=1)
	return rank_1_freq, scores
\end{verbatim}

For TVK, the simulation randomly samples admissible $N$-models within perturbed ratio intervals and tests component-wise dominance on expanded vectors.

\subsubsection{Practical Recommendations}

Based on the robustness analysis:

\begin{enumerate}
	\item \textbf{Use SMART} when weights are uncertain but equal treatment of criteria is acceptable; it provides the most stable rankings under perturbation.
	
	\item \textbf{Use SMARTS} when swing-weighting information is available; it balances interpretability with moderate robustness.
	
	\item \textbf{Use TVK} when preference information is inherently interval-based or criteria are ordinal; accept that some alternatives may be declared incomparable under high uncertainty.
	
	\item \textbf{Hybrid approach}: Run all three methods and compare:
	\begin{itemize}
		\item If all agree on the top alternative $\rightarrow$ high confidence recommendation
		\item If TVK declares incomparability while SMART/SMARTS rank $\rightarrow$ investigate sensitivity to importance ratios
		\item If rankings diverge significantly $\rightarrow$ collect more precise preference information
	\end{itemize}
\end{enumerate}

\begin{thebibliography}{99}
	
	\bibitem{edwards1994smarts}
	W. Edwards and F.H. Barron, ``SMARTS and SMARTER: Improved simple methods for multiattribute utility measurement,'' \emph{Organ. Behav. Hum. Decis. Process.}, vol. 60, no. 1, pp. 26--37, 1994.
	
	\bibitem{podinovskii2019idei}
	V.V. Podinovskii, \emph{Ideas and Methods of Criteria Importance Theory in Multi-Criteria Decision Problems}. Moscow: Nauka, 2019.
	
	\bibitem{podinovskii2022mcdm}
	V.V. Podinovskii, \emph{Metody mnogokriterial'nogo analiza reshenij}. Moscow: Fizmatlit, 2022.
	
	\bibitem{preflib2026}
	PrefLib Team, ``Format specification and dataset repository,'' \href{https://preflib.github.io/PrefLib-Jekyll/format}{preflib.github.io/PrefLib-Jekyll/format}
	
	\bibitem{mao2013better}
	A. Mao, A.D. Procaccia, and Y. Chen, ``Better human computation through principled voting,'' in \emph{Proc. AAAI Conf. Artif. Intell.}, vol. 27, no. 1, pp. 1124--1130, 2013.
	
	\bibitem{petrovsky2017}
	A.B. Petrovsky, \emph{Teoriya resheniy: Vybor na mnozhestve kriteriev}. Moscow: LENAND, 2017.
	
	\bibitem{brams2007approval}
	S.J. Brams and P.C. Fishburn, \emph{Approval Voting}, 2nd ed. Springer, 2007.
	
	\bibitem{simpson1969defining}
	P.B. Simpson, ``Defining a representation independent degree of distortion for voting systems,'' \emph{Behav. Sci.}, vol. 14, pp. 41--52, 1969.
	
	\bibitem{dodgson1876method}
	C.L. Dodgson, \emph{A Method of Taking Votes on More than Two Issues}. Oxford: Clarendon Press, 1876.
	
	\bibitem{caragiannis2016dodgson}
	I. Caragiannis et al., ``On the approximability of Dodgson and Young elections,'' \emph{Artif. Intell.}, vol. 231, pp. 19--43, 2016.
	
	\bibitem{bartholdi1989voting}
	J.J. Bartholdi, C.A. Tovey, and M.A. Trick, ``Voting schemes for which it can be difficult to tell who won the election,'' \emph{Soc. Choice Welf.}, vol. 6, no. 2, pp. 157--165, 1989.
	
	\bibitem{parkhomenko2024pydass}
	V.A. Parkhomenko and M.A. Lopatin, ``Method of sequential preference refinement,'' Lecture Notes, SPbPU, St. Petersburg, 2024.
	
\end{thebibliography}
\section*{Availability of Code and Data}
The complete source code for this project, including all MCDM method implementations and case study data, is available on GitHub:
\begin{center}
	\url{https://github.com/Majd-nakhla/report-Comparative-Analysis-of-SMART-SMARTS-and-TVK}
\end{center}

\end{document}
